\section{Thoughts around the Maypole}

\begin{quotex}
When we plunge into those mute and motionless depths in which the muddy stream of our actuality has its source, being careful not to violate its purity and peace, we come into inward contact, in this source of our own spiritual life, with the source of the universal life.
\flright{\textsc{Vladimir Solovyov}}
\end{quotex}


While some people can only see a certain world leader as a thug, we are incredulous that a mere thug would be reading and recommending thinkers like \textbf{Vladimir Solovyov} or \textbf{Nicolai Berdyaev}. Perhaps we will go through the former's Justification of the Good to look for clues to world events. For now, however, we are still trying to explicate Solovyov's description of creation through God's three positings. It should be read in conjunction with \textbf{Rene Guenon}'s Multiple States of Being, as it fills in the missing gap from possibility to actuality.

\paragraph{The Artist and the Creative Act}


Solovyov uses the example of the creativity of the artist as analogy, another example of the Hermetic method. It is worth meditating on this. Of the forms of man's spiritual activity in the world -- religion, art, politics, science, economics -- why does Solovyov choose the artist? Like God, the artist holds his artistic ideas in his consciousness. Solovyov explains:

\begin{quotex}
The artistic idea is certainly not alien or external to the artist. It is the artist's own inner essence, the core of his spirit and the content of his life; it is what makes him what he is. In striving to realize or embody this idea in an actual work of art, the artist wishes merely to have this inner essence, this idea, not only within himself, but also for himself, or before himself as an object. The artist wishes to represent his own as ``another'', in another, objective form.
\end{quotex}


This applies especially to the writer, who strives to objectify his thoughts in written form. This form should be beautiful, coherent, intelligible, accurate, grammatical, rhetorically stimulating, and so on. Solovyov does qualify this analogy when he admits that

\begin{quotex}
Artistic creation presupposes a certain passive state of inspiration or inner reception, in which the artist does not possess an idea but is possessed by it.
\end{quotex}


This just means that the idea is prior to any manifestation of it and arises out of non-being. Even as the idea is being formed as a thought in consciousness, that is the beginning of objectification.
\textbf{Valentin Tomberg} develops this concept fully in the fourth Letter, as he discusses the stages from inspiration to the ``book''.

So the first positing is the idea in mind, and the second is its objectification. This object is both other and not-other. This unity of the object with the creator is the third positing. That is, the object is part of multiplicity, yet it is a manifestation of the unity in the artist's consciousness.

Is the analogy clear? Or is another inspiration necessary?

\paragraph{Microcosms and Heaps}


Organic unity is the expression of one spirit in multiplicity. A heap is collection of independent units with little more than proximity in common. An organism, therefore, is the objectification of one idea. A heap is the manifestation of many ideas, even if they seem to be somewhat related.

The modern mind cannot focus, has no power of concentration, and is therefore attracted to heaps. It loves the intellectual diversity, just as some people prefer the variety at the food court in their local mall. There is an intellectual thrill that comes with hearing new ideas, very like the feeling of a boy opening his birthday gifts. As such it is part of the glamour of the world, a tinsel attraction, and just adds to the chaos in the noosphere.

Those organizations that pretend to oppose the modern world need to be true to the etymology of the word, i.e., it should be an organic unity, rather than a heap. The organization, the conference, needs to be a microcosm, a replication of the new society it has in mind. This is a fourfold vision. Otherwise, it devolves into a single vision, Newton's sleep, a simple concept that is the least common denominator. There is multiplicity with no unity.

The political as artistic creation.


\flrightit{Posted on 2014-05-01 by Cologero}

